\section{Numerical implementation and verification}
\label{sec:events}\label{sec:numerics}\label{sec:events:exact}
\label{sec:numerics:fv}\label{sec:numerics:reject}\label{sec:numerics:acceptance}
This section explains how the particle equations are integrated in time
and how accepted states are checked for admissibility and conservation.
The numerical calculation advances component inventories and energy under
the prescribed gas state, heat supply and exposure. Surface composition,
particle temperature and front motion are solved responses. Each phase
regime uses the corresponding equations of Section~\ref{sec:model}; the
numerical scheme transfers the conserved quantities across regime changes.

\subsection{Spatial discretization and coupled time integration}
Regional balances are integrated in spherical finite volumes. In regime A,
the present experiments represent the particle as a uniform thermal body,
coupled to the liquid surface through a one-cell conductive resistance.
This reduces the coupled startup problem to one body temperature, while
retaining a separate surface temperature and conductive heat transfer.
It is a coarse approximation, not a spatially resolved preheating
calculation; its effect on the initial temperature profile should be
assessed before detailed startup heating is interpreted. Regimes B and C
resolve the radial particle.
The receding interface cuts a fixed radial grid and is represented by the
wet volume fraction $\frontz=(\front/\Rp)^3$. Its swept-volume rate is
\begin{equation}
Q_{\Gamma}=V_{\Rp}\frac{\frontz_{\mathrm{new}}-\frontz_{\mathrm{old}}}{\Delta t}.
\label{eq:sweptvolume}
\end{equation}
The same rate enters both regional balances and the interfacial jumps.
Backward-Euler integration solves temperatures, compositions, external
liquid inventories and front position together. Component enthalpy transport
and conductive heat use the same face states as the storage calculation.
Detailed face reconstruction and nonlinear controls are given in
Sec.~\ref*{S-sup:activation}.

\subsection{Phase changes and moving-interface events}
When external free hexane is depleted, the dry shell can appear from zero
thickness.
The first finite-volume interval satisfies the component and energy jumps
without assigning an arbitrary initial shell. Arrival at and departure
from radial grid faces are solved as geometric events. Near the centre,
small positive-core intervals precede a conservative projection to the
zero-core particle. That projection includes the external exchanges over
its estimated duration; the conditional remaining-time allowance is below
2 ms for DT and 20 ms for Faner. These bounds concern the terminal event,
rather than the error of the complete trajectory.
External water appearance and depletion likewise switch the surface
equations while preserving water and energy. The liquid-water imbibition pathway
operates only in a dry shell connected to external liquid. The DT
trajectory retains positive external water throughout regime B, so this
experiment does not exercise the moving front without a water film.
Regime C includes liquid-present, contact and liquid-absent surface states.

\subsection{Physical admissibility and conservation checks}
At each time step, the nonlinear solution must satisfy non-negative
inventories, pressure-feasible
phase equilibrium, the receding-front geometry and the retained-water
capacity. Capacity is expressed as an active set,
\begin{equation}
0\le\Wcap-\Wret\;\perp\;\lambda\ge0,
\label{eq:capcomplementarity}
\end{equation}
with $\lambda$ its dimensionless dual variable. Local component, energy and
combined inventory balances are checked at a normalized tolerance of
$10^{-10}$. Independent accumulation of the external exchanges provides
a separate whole-history check. The four-cell Faner reconstruction uses a
reporting allowance of $4\times10^{-10}$, corresponding to at most
$10^{-11}$ kg water/kg dry meal; its step acceptance tolerance
remains $10^{-10}$.

For regular Faner steps, the local nonlinear root criterion can be
$10^{-8}$ while the component and energy acceptance checks retain
$10^{-10}$. Near wet volume fractions below $10^{-7}$, the root criterion
can be $2\times10^{-8}$ and the scaled conditioning ceiling $10^{18}$.
These numerical controls do not change the component or energy equations,
their conservation tolerance, or phase admissibility. The Supplementary
Material reports the root residual, the conserved-inventory error and the
terminal-time estimate separately. Halving the stationary-particle
timestep and comparing two and four radial cells quantify particular
sensitivities, not whole-history mesh convergence
(Section~\ref{sec:discussion:limits}).
